\chapter{Sistema de coordenadas y segmentos en el plano}

\begin{center}
\begin{minipage}{.88\textwidth}
\small\sffamily
\textbf{Propósito.} Ubicar puntos y resolver problemas geométricos mediante
distancia, punto medio y división de segmentos.
\end{minipage}
\end{center}

\section{El plano cartesiano}

\begin{definition}{Plano cartesiano}{plano}
El plano cartesiano está formado por dos rectas numéricas perpendiculares: el
eje horizontal $x$ y el eje vertical $y$. Su intersección es el origen
$O(0,0)$. Un punto se representa mediante un par ordenado $P(x,y)$.
\end{definition}

El signo de las coordenadas determina el cuadrante:
\[
\text{I }(+,+),\qquad \text{II }(-,+),\qquad
\text{III }(-,-),\qquad \text{IV }(+,-).
\]
Los puntos sobre los ejes no pertenecen a ningún cuadrante.

\begin{example}{Ubicación de puntos}{ubicar}
Ubica $A(3,2)$, $B(-2,3)$, $C(-3,-2)$ y $D(2,-3)$.

\begin{center}
\begin{tikzpicture}
\begin{axis}[ga axis,width=8cm,height=7cm,xmin=-4,xmax=4,ymin=-4,ymax=4,
  xtick={-3,...,3},ytick={-3,...,3}]
\addplot[only marks,mark=*,mark size=2.2pt,GArosa]
  coordinates {(3,2) (-2,3) (-3,-2) (2,-3)};
\node[above right] at (axis cs:3,2) {$A$};
\node[above left] at (axis cs:-2,3) {$B$};
\node[below left] at (axis cs:-3,-2) {$C$};
\node[below right] at (axis cs:2,-3) {$D$};
\end{axis}
\end{tikzpicture}
\end{center}
Los puntos están, respectivamente, en los cuadrantes I, II, III y IV.
\end{example}

\begin{exercise}{Coordenadas y cuadrantes}{cuadrantes}
Indica el cuadrante o eje al que pertenece cada punto:
$P(-5,2)$, $Q(0,-4)$, $R(3,-1)$, $S(-2,-6)$ y $T(4,0)$.
\end{exercise}

\section{Distancia entre dos puntos}

Sean $P_1(x_1,y_1)$ y $P_2(x_2,y_2)$. Las diferencias horizontal y vertical
son
\[
\Delta x=x_2-x_1,\qquad \Delta y=y_2-y_1.
\]

\begin{intuition}
Estas diferencias son los catetos de un triángulo rectángulo. La distancia
entre los puntos es su hipotenusa.
\end{intuition}

Por el teorema de Pitágoras,
\[
d^2=(\Delta x)^2+(\Delta y)^2.
\]

\begin{formula}
\[
\boxed{d(P_1,P_2)=\sqrt{(x_2-x_1)^2+(y_2-y_1)^2}}
\]
\end{formula}

\begin{example}{Cálculo de una distancia}{distancia}
Calcula la distancia entre $A(-2,1)$ y $B(4,5)$.
\begin{solution}
\[
d=\sqrt{(4-(-2))^2+(5-1)^2}
 =\sqrt{6^2+4^2}=\sqrt{52}=2\sqrt{13}.
\]
La distancia es $2\sqrt{13}\approx 7.21$ unidades.
\end{solution}
\end{example}

\begin{remark}
El orden de los puntos no modifica la distancia porque las diferencias están
elevadas al cuadrado.
\end{remark}

\begin{exercise}{Práctica de distancia}{practica-distancia}
Calcula la distancia entre cada pareja.
\begin{enumerate}[label=\alph*)]
  \item $A(1,2)$ y $B(5,5)$.
  \item $C(-3,4)$ y $D(2,-8)$.
  \item $E(-4,-1)$ y $F(-4,6)$.
  \item $G(2,-3)$ y el origen.
\end{enumerate}
\end{exercise}

\section{Punto medio de un segmento}

El punto medio se obtiene promediando por separado las coordenadas de los
extremos.

\begin{formula}
Si $A(x_1,y_1)$ y $B(x_2,y_2)$, entonces
\[
\boxed{M\left(\frac{x_1+x_2}{2},\frac{y_1+y_2}{2}\right)}.
\]
\end{formula}

\begin{example}{Punto medio}{medio}
Halla el punto medio del segmento cuyos extremos son $A(-5,3)$ y $B(7,-1)$.
\begin{solution}
\[
M=\left(\frac{-5+7}{2},\frac{3+(-1)}{2}\right)=(1,1).
\]
\end{solution}
\end{example}

\begin{example}{Extremo desconocido}{extremo}
El punto medio de $AB$ es $M(3,2)$ y $A(-1,5)$. Determina $B(x,y)$.
\begin{solution}
\[
3=\frac{-1+x}{2}\Rightarrow x=7,
\qquad
2=\frac{5+y}{2}\Rightarrow y=-1.
\]
Por tanto, $B(7,-1)$.
\end{solution}
\end{example}

\begin{exercise}{Puntos medios}{practica-medio}
\begin{enumerate}[label=\alph*)]
  \item Halla el punto medio de $A(2,-3)$ y $B(8,5)$.
  \item Halla el punto medio de $C(-7,4)$ y $D(3,-6)$.
  \item Si $M(-2,3)$ es el punto medio de $A(4,1)$ y $B$, determina $B$.
\end{enumerate}
\end{exercise}

\section{División de un segmento en una razón dada}

Sea $P$ un punto que divide internamente el segmento $AB$ en la razón
\[
AP:PB=m:n.
\]
Entonces $P$ está más cerca del extremo asociado con el peso menor.

\begin{formula}
Para $A(x_1,y_1)$ y $B(x_2,y_2)$,
\[
\boxed{P\left(\frac{nx_1+mx_2}{m+n},
                \frac{ny_1+my_2}{m+n}\right)}.
\]
\end{formula}

\begin{remark}
La coordenada de $A$ se multiplica por $n$ y la de $B$ por $m$: los pesos
aparecen cruzados. Si $m=n$, la fórmula se reduce al punto medio.
\end{remark}

\begin{example}{División interna}{division}
El punto $P$ divide el segmento de $A(1,2)$ a $B(10,8)$ en la razón $AP:PB=2:1$.
\begin{solution}
Aquí $m=2$ y $n=1$:
\[
P=\left(\frac{1(1)+2(10)}{3},
         \frac{1(2)+2(8)}{3}\right)=(7,6).
\]
Como $AP$ es dos veces $PB$, $P$ queda más cerca de $B$.
\end{solution}
\end{example}

\begin{exercise}{División de segmentos}{practica-division}
\begin{enumerate}[label=\alph*)]
  \item Divide el segmento de $A(-2,1)$ a $B(7,4)$ en razón $1:2$.
  \item Divide el segmento de $C(3,-5)$ a $D(-7,5)$ en razón $3:2$.
  \item Comprueba con la fórmula que la razón $1:1$ produce el punto medio.
\end{enumerate}
\end{exercise}

\section{Aplicaciones geométricas}

\subsection{Perímetro y clasificación de triángulos}

Para hallar un perímetro se calcula la longitud de cada lado y se suman los
resultados. Comparar las longitudes permite clasificar un triángulo.

\begin{example}{Triángulo en el plano}{triangulo}
Sean $A(0,0)$, $B(4,0)$ y $C(0,3)$.
\begin{solution}
\[
AB=4,\qquad AC=3,\qquad BC=\sqrt{(-4)^2+3^2}=5.
\]
El perímetro es $3+4+5=12$. Además,
\[
3^2+4^2=5^2,
\]
así que el triángulo es rectángulo.
\end{solution}
\end{example}

\subsection{Área de un triángulo}

Para $A(x_1,y_1)$, $B(x_2,y_2)$ y $C(x_3,y_3)$:
\begin{formula}
\[
\boxed{\mathcal A=\frac12\left|
x_1(y_2-y_3)+x_2(y_3-y_1)+x_3(y_1-y_2)
\right|}.
\]
\end{formula}

\begin{example}{Área mediante coordenadas}{area}
Calcula el área del triángulo $A(1,1)$, $B(5,1)$ y $C(3,4)$.
\begin{solution}
\[
\mathcal A=\frac12|1(1-4)+5(4-1)+3(1-1)|
=\frac12|-3+15|=6.
\]
El área es $6$ unidades cuadradas.
\end{solution}
\end{example}

\section{Ejercicios de unidad}

\begin{exercise}{Ejercicios básicos}{unidad-basicos}
\begin{enumerate}
  \item Grafica $A(-4,3)$, $B(2,5)$, $C(3,-2)$ y $D(-1,-4)$.
  \item Calcula la distancia entre $P(-6,2)$ y $Q(2,-4)$.
  \item Halla el punto medio de $R(7,-3)$ y $S(-5,9)$.
  \item El punto medio de $AB$ es $M(2,-1)$. Si $A(-3,4)$, halla $B$.
  \item Divide el segmento de $A(0,0)$ a $B(12,9)$ en la razón $2:1$.
  \item Verifica si $A(-1,1)$, $B(3,1)$ y $C(1,1+2\sqrt3)$ forman un
  triángulo equilátero.
\end{enumerate}
\end{exercise}

\begin{problem}{Aplicaciones}{unidad-aplicaciones}
\begin{enumerate}
  \item Los vértices de un terreno son $A(1,1)$, $B(7,1)$, $C(7,5)$ y
  $D(1,5)$. Calcula su perímetro y área.
  \item Determina si el triángulo $A(-2,1)$, $B(4,3)$ y $C(2,-3)$ es
  isósceles. Justifica mediante distancias.
  \item Un poste se colocará a la tercera parte del trayecto de $A(-3,2)$ a
  $B(9,8)$, medida desde $A$. Obtén sus coordenadas.
  \item Calcula el área del triángulo con vértices $(-2,-1)$, $(4,2)$ y
  $(1,6)$.
  \item Demuestra mediante distancias que los puntos $A(0,0)$, $B(6,2)$,
  $C(4,8)$ y $D(-2,6)$ forman un cuadrado.
\end{enumerate}
\end{problem}

\begin{hints}
En el último problema compara los cuatro lados y las dos diagonales. Un
cuadrado tiene cuatro lados iguales y diagonales iguales.
\end{hints}

\enlargethispage{5\baselineskip}
\section*{Resumen}
\addcontentsline{toc}{section}{Resumen}

\begin{properties}
Para $A(x_1,y_1)$ y $B(x_2,y_2)$:
\[
\begin{aligned}
d(A,B)&=\sqrt{(x_2-x_1)^2+(y_2-y_1)^2},\\[2mm]
M&=\left(\frac{x_1+x_2}{2},\frac{y_1+y_2}{2}\right),\\[2mm]
P_{m:n}&=\left(\frac{nx_1+mx_2}{m+n},
\frac{ny_1+my_2}{m+n}\right).
\end{aligned}
\]
\end{properties}
